\begin{table}[p]
\centering
\sffamily
\scriptsize
\setlength{\tabcolsep}{4.4pt}
\renewcommand{\arraystretch}{1.03}
\begin{threeparttable}
\caption{Niveaux relatifs et rangs dans l'OCDE en 2019 (États-Unis = 100)}
\label{tab:note_lp_levels}
\begin{tabular}{l *{4}{rr}}
\toprule
& \multicolumn{2}{c}{$Y/L$} & \multicolumn{2}{c}{PTF} & \multicolumn{2}{c}{$K/Y$} & \multicolumn{2}{c}{$K/L$} \\
\cmidrule(lr){2-3} \cmidrule(lr){4-5} \cmidrule(lr){6-7} \cmidrule(lr){8-9}
& Ratio & Rang & Ratio & Rang & Ratio & Rang & Ratio & Rang \\
\midrule
Irlande & 179 & 1 & 94 & 10 & 151 & 9 & 270 & 1 \\
Norvège & 148 & 2 & 110 & 1 & 123 & 23 & 181 & 2 \\
Suisse & 119 & 3 & 97 & 7 & 150 & 10 & 178 & 3 \\
Luxembourg & 115 & 4 & 101 & 3 & 114 & 27 & 131 & 9 \\
Danemark & 109 & 5 & 94 & 9 & 140 & 14 & 152 & 7 \\
Autriche & 104 & 6 & 90 & 14 & 153 & 8 & 159 & 5 \\
Allemagne & 103 & 7 & 104 & 2 & 105 & 28 & 109 & 16 \\
France & 103 & 8 & 94 & 11 & 153 & 7 & 157 & 6 \\
Suède & 102 & 9 & 91 & 12 & 146 & 13 & 149 & 8 \\
Belgique & 101 & 10 & 90 & 13 & 159 & 4 & 160 & 4 \\
Pays-Bas & 101 & 11 & 98 & 5 & 119 & 26 & 120 & 13 \\
États-Unis & 100 & 12 & 100 & 4 & 100 & 29 & 100 & 18 \\
Australie & 90 & 13 & 85 & 19 & 131 & 21 & 118 & 14 \\
Finlande & 89 & 14 & 83 & 21 & 139 & 15 & 124 & 11 \\
Islande & 88 & 15 & 90 & 15 & 138 & 16 & 122 & 12 \\
Royaume-Uni & 85 & 16 & 84 & 20 & 121 & 25 & 102 & 17 \\
Italie & 84 & 17 & 79 & 24 & 154 & 6 & 130 & 10 \\
Espagne & 80 & 18 & 87 & 17 & 137 & 17 & 110 & 15 \\
\textbf{Canada} & \textbf{77} & \textbf{19} & \textbf{82} & \textbf{22} & \textbf{121} & \textbf{24} & \textbf{94} & \textbf{20} \\
Tchéquie & 63 & 20 & 64 & 34 & 147 & 12 & 93 & 21 \\
Corée du Sud & 63 & 21 & 67 & 31 & 136 & 18 & 86 & 24 \\
Slovénie & 63 & 22 & 67 & 30 & 156 & 5 & 99 & 19 \\
Turquie & 63 & 23 & 98 & 6 & 86 & 33 & 54 & 31 \\
Nouvelle-Zélande & 63 & 24 & 89 & 16 & 82 & 34 & 52 & 33 \\
Japon & 62 & 25 & 66 & 32 & 148 & 11 & 92 & 22 \\
Estonie & 60 & 26 & 71 & 27 & 125 & 22 & 75 & 26 \\
Pologne & 58 & 27 & 97 & 8 & 63 & 37 & 37 & 34 \\
Lituanie & 57 & 28 & 79 & 25 & 97 & 31 & 55 & 30 \\
Israël & 54 & 29 & 71 & 28 & 98 & 30 & 53 & 32 \\
Portugal & 54 & 30 & 72 & 26 & 161 & 3 & 87 & 23 \\
Hongrie & 52 & 31 & 64 & 33 & 132 & 20 & 68 & 28 \\
Lettonie & 51 & 32 & 62 & 35 & 165 & 2 & 84 & 25 \\
Slovaquie & 47 & 33 & 58 & 37 & 133 & 19 & 63 & 29 \\
Grèce & 45 & 34 & 57 & 38 & 165 & 1 & 75 & 27 \\
Chili & 41 & 35 & 87 & 18 & 59 & 38 & 24 & 36 \\
Mexique & 37 & 36 & 68 & 29 & 96 & 32 & 35 & 35 \\
Costa Rica & 33 & 37 & 79 & 23 & 66 & 36 & 22 & 37 \\
Colombie & 20 & 38 & 59 & 36 & 77 & 35 & 15 & 38 \\
\bottomrule
\end{tabular}
\begin{tablenotes}\footnotesize
\item \textit{Note}\,: Ratios relatifs aux États-Unis ($=100$) et rangs calculés sur 38 membres de l'OCDE. Les colonnes sont comparables comme ratios relatifs, mais elles ne correspondent pas au même objet économique\,: $Y/L$ et $K/L$ sont des niveaux par heure, la PTF est un indice de niveau, et $K/Y$ un ratio capital/production. Voir l'annexe pour les définitions et les détails de construction. Source\,: Penn World Table 11.0.
\end{tablenotes}
\end{threeparttable}
\end{table}
